\documentclass[11pt]{article}

\usepackage[T1]{fontenc}
\usepackage[utf8]{inputenc}
\usepackage{lmodern}
\usepackage{amsmath,amssymb,amsthm,mathtools}
\usepackage{geometry}
\usepackage{booktabs}
\usepackage{longtable}
\usepackage{array}
\usepackage{hyperref}
\usepackage{enumitem}

\geometry{margin=1in}
\setlist[itemize]{leftmargin=1.5em}
\setlist[enumerate]{leftmargin=1.75em}

\newtheorem{remark}{Remark}
\newtheorem{observation}{Observation}
\newcommand{\code}[1]{\text{\texttt{#1}}}

\title{A Formula Dictionary for the HAETAE NTT Verification in EasyCrypt}
\author{Codex Proof Note}
\date{22 April 2026}

\begin{document}
\maketitle

\begin{abstract}
This note organizes the principal formulas underlying the HAETAE NTT
verification project in EasyCrypt. Its purpose is not to restate the current
proof status, but rather to index the development by mathematical task: for
each task in the project, we record the central identities, invariants,
representation maps, and target equivalences that must be established. The
result is a formula-oriented guide to the project, written in the style of a
mathematical aide-memoire for cryptographic verification.
\end{abstract}

\section{Global Verification Objective}

Let \(p\) be a polynomial represented concretely by a byte array
\(\texttt{rp} : \texttt{BArray1024.t}\). The final end-to-end objective of the
project is the pair of Hoare-style statements
\[
\{\code{poly\_repr}(\texttt{rp},p)\}
\ \texttt{poly\_ntt\_jazz}(\texttt{rp})\
\{\code{poly\_repr}(\texttt{res}, \mathrm{Rq.ntt}(p))\},
\]
\[
\{\code{poly\_repr}(\texttt{rp},p)\}
\ \texttt{poly\_invntt\_jazz}(\texttt{rp})\
\{\code{poly\_repr}(\texttt{res}, \mathrm{Rq.invntt}(p))\}.
\]

Every other formula in the project serves this final pair of correctness
theorems.

\section{Ambient Notation}

The field is \(\mathbf{F}_q\) with
\[
q = 64513.
\]
The transform length is
\[
N = 256,
\]
and the chosen primitive \(512\)-th root of unity is
\[
zroot = 426 \in \mathbf{F}_q.
\]
The Montgomery constant used throughout the extracted arithmetic is
\[
R = 2^{32} \bmod q = 14321.
\]
The bit-reversal permutation on \(8\) bits is written
\[
\operatorname{br}(i) = \operatorname{bsrev}_8(i).
\]

\section{Task I: Imperative NTT Specifications}

The forward imperative specification \(\texttt{NTT.ntt}\) executes
Cooley--Tukey butterflies with layer lengths
\[
128, 64, 32, 16, 8, 4, 2, 1.
\]
The inverse imperative specification \(\texttt{NTT.invntt}\) executes inverse
layers with lengths
\[
1, 2, 4, 8, 16, 32, 64, 128,
\]
followed by a global scaling by \(\operatorname{inv}(256)\).

The forward butterfly is
\[
t = \zeta \cdot r_{j+\ell},\qquad
r_{j+\ell}' = r_j - t,\qquad
r_j' = r_j + t.
\]
The inverse butterfly is
\[
t = r_j,\qquad
r_j' = t + r_{j+\ell},\qquad
r_{j+\ell}' = \zeta \cdot (t - r_{j+\ell}).
\]

The formal tasks attached to these formulas are:
\begin{itemize}
\item prove losslessness of \(\texttt{NTT.ntt}\);
\item prove losslessness of \(\texttt{NTT.invntt}\);
\item use these imperative models as refinement targets for the extracted code.
\end{itemize}

\section{Task II: Termination Variants}

For the forward procedure the variants are:
\[
V_{\mathrm{outer}} = \texttt{len},\qquad
V_{\mathrm{middle}} = 256 - \texttt{start},\qquad
V_{\mathrm{inner}} = \texttt{start} + \texttt{len} - \texttt{j}.
\]
For the inverse procedure one uses:
\[
V_{\mathrm{outer}} = 256 - \texttt{len},\qquad
V_{\mathrm{middle}} = 256 - \texttt{start},\qquad
V_{\mathrm{inner}} = \texttt{start} + \texttt{len} - \texttt{j},
\]
and for the final scaling loop
\[
V_{\mathrm{scale}} = 256 - \texttt{j}.
\]

These formulas belong to the purely structural task of proving
\texttt{islossless}. They do not encode semantics of the transform itself, but
they are indispensable for any subsequent program logic argument.

\section{Task III: Root-of-Unity Structure}

The transform depends on the order of \(zroot\). The essential identities are
\[
zroot^{128} = 28837,
\qquad
zroot^{256} = -1,
\qquad
zroot^{512} = 1.
\]
In EasyCrypt, the decisive lemma is
\[
\texttt{Zq.exp zroot 256 = incoeff (-1)}.
\]

This task supports:
\begin{itemize}
\item the correctness of the inverse transform;
\item the algebraic interpretation of the twiddle tables;
\item the final normalization by \(\operatorname{inv}(256)\).
\end{itemize}

\section{Task IV: Abstract NTT Formulas}

The abstract forward transform in \texttt{Rq.ec} is
\[
\mathrm{ntt}(p)[i]
=
\begin{cases}
\displaystyle \sum_{j=0}^{127} p_{2j}\, zroot^{(2\operatorname{br}(i/2)+1)j},
& i \equiv 0 \pmod 2,\\[1.25em]
\displaystyle \sum_{j=0}^{127} p_{2j+1}\, zroot^{(2\operatorname{br}(i/2)+1)j},
& i \equiv 1 \pmod 2.
\end{cases}
\]

The abstract inverse transform is
\[
\mathrm{invntt}(p)[i]
=
\begin{cases}
\displaystyle \sum_{j=0}^{127} \operatorname{inv}(128)\, p_{2j}\,
zroot^{-(2\operatorname{br}(j)+1)(i/2)},
& i \equiv 0 \pmod 2,\\[1.25em]
\displaystyle \sum_{j=0}^{127} \operatorname{inv}(128)\, p_{2j+1}\,
zroot^{-(2\operatorname{br}(j)+1)(i/2)},
& i \equiv 1 \pmod 2.
\end{cases}
\]

This is the target mathematical semantics to which the imperative routines
must eventually be shown equal.

\section{Task V: Representation of Concrete Program State}

The concrete-to-abstract bridge is given by
\[
\code{word\_to\_coeff}(w) = \operatorname{incoeff}(\code{to\_sint}(w)),
\]
\[
\code{barray256\_to\_poly}(\texttt{rp})[i]
=
\code{word\_to\_coeff}(\texttt{BArray1024.get32}\; \texttt{rp}\; i),
\]
and
\[
\code{poly\_repr}(\texttt{rp},p)
:\Longleftrightarrow
\code{barray256\_to\_poly}(\texttt{rp}) = p.
\]

The formulas attached to this task are:
\begin{itemize}
\item pointwise readback:
  \[
  \code{poly\_repr}(\texttt{rp},p) \Longrightarrow
  p[i] = \code{word\_to\_coeff}(\texttt{BArray1024.get32}\; \texttt{rp}\; i);
  \]
\item small signed-word conversion:
  \[
  |z| < q \Longrightarrow
  \code{word\_to\_coeff}(\texttt{W32.of\_int}\; z) = \operatorname{incoeff}(z);
  \]
\item normalized representative conversion:
  \[
  |z| < q \Longrightarrow
  \code{word\_to\_coeff}(\texttt{W32.of\_int}\; z)
  =
  \operatorname{incoeff}\!\bigl(\mathbf{if}\ 0 \le z\ \mathbf{then}\ z\ \mathbf{else}\ z+q\bigr).
  \]
\end{itemize}

This task is the basic language in which all simulation statements are
written.

\section{Task VI: Twiddle Tables}

\subsection{Forward Table}

The abstract forward twiddle table is
\[
\zeta_k = zroot^{\operatorname{br}(k)}.
\]
Its Montgomery form is
\[
\zeta_k^{\mathrm{Mont}} = \zeta_k \cdot R.
\]
The concrete table \texttt{jzetas} stores precisely these Montgomery values,
except that index \(0\) is padding:
\[
\texttt{jzetas}[0] = 0,
\qquad
\texttt{jzetas}[k] \leftrightsquigarrow \zeta_k \cdot R
\quad (1 \le k \le 255).
\]

Thus the intended bridge formula is
\[
\code{barray256\_to\_poly}(\texttt{jzetas})
=
(\code{array256\_mont}(\zeta))[0 \leftarrow 0].
\]

\subsection{Inverse Table}

The abstract inverse twiddle table is
\[
\zeta_k^{-1} = zroot^{-(\operatorname{br}(k)+1)}
\quad (0 \le k \le 254),
\]
while the last slot is the inverse normalization:
\[
\zeta_{255}^{-1} = \operatorname{inv}(256).
\]
Its Montgomery form is
\[
(\zeta_k^{-1})^{\mathrm{Mont}} = \zeta_k^{-1}\cdot R.
\]
The concrete table \texttt{jzetas\_inv} matches this at indices
\[
0 \le k \le 254,
\]
but at index \(255\) it stores the machine value
\[
-29720,
\]
which encodes the final scaling constant in the format expected by the
Montgomery multiplication routine.

Accordingly, the bridge formula is
\[
\code{barray256\_to\_poly}(\texttt{jzetas\_inv})
=
(\code{array256\_mont\_inv}(\zeta^{-1}))[255 \leftarrow -29720].
\]

\begin{remark}
This is the most delicate formula in the current project state. It is not
conceptually difficult, but it sits precisely at the interface between the
mathematics of twiddle factors and the byte-level normalization performed by
the EasyCrypt libraries.
\end{remark}

\section{Task VII: Inverse Scaling Constant}

Two related constants appear in the inverse transform:
\[
\code{scale255} = \operatorname{inv}(256) = -252 \pmod q,
\]
and the machine value
\[
f = -29720.
\]
The key identity is the Montgomery compatibility relation
\[
\operatorname{incoeff}(f)\cdot R^{-1}
=
\code{scale255}\cdot R.
\]

This formula serves the final loop of the inverse extracted code, where the
program multiplies each coefficient by the machine word stored at
\texttt{jzetas\_inv[255]} via the Montgomery multiplication routine.

\section{Task VIII: Montgomery Multiplication}

The extracted routine \texttt{M.\_\_fqmul} computes
\[
\mathrm{MontMul}(a,b)
=
a\cdot b\cdot R^{-1} \pmod q,
\]
with the result again in Montgomery representation when both inputs are
already Montgomery-encoded.

The intended correctness formula is
\[
\code{word\_to\_coeff}(\texttt{M.\_\_fqmul}(a,b))
=
\code{word\_to\_coeff}(a)\cdot \code{word\_to\_coeff}(b)\cdot R.
\]

This task depends on the Montgomery-reduction theorem from
\texttt{Montgomery.ec} and supplies the arithmetic core of every butterfly
simulation.

\section{Task IX: Forward Butterfly Simulation}

Let the concrete word read from the table represent \(\zeta\cdot R\). Let the
current concrete coefficients represent \(x\) and \(y\). Then the extracted
forward butterfly computes, at the field level,
\[
t = \zeta y,
\qquad
x' = x + t,
\qquad
y' = x - t.
\]

Thus the local proof obligation is:
\[
\text{concrete forward butterfly}
\Longrightarrow
\text{imperative forward butterfly in } \texttt{NTT.ntt}.
\]

This task combines the formulas from:
\begin{itemize}
\item the representation map;
\item the forward twiddle bridge;
\item Montgomery multiplication correctness;
\item field addition and subtraction correctness.
\end{itemize}

\section{Task X: Inverse Butterfly Simulation}

For the inverse transform one must prove the corresponding local statement:
\[
t = x,
\qquad
x' = x + y,
\qquad
y' = \zeta(x - y).
\]

At the program level the order of updates matters, so the precise proof must
track the concrete sequence:
\[
t \leftarrow r_j,\quad
r_j \leftarrow t + r_{j+\ell},\quad
r_{j+\ell} \leftarrow t - r_{j+\ell},\quad
r_{j+\ell} \leftarrow \zeta \cdot r_{j+\ell}.
\]

This task is the inverse analogue of the previous section, augmented by the
final scaling formula from Task VII.

\section{Task XI: Loop Invariants for the Extracted Jasmin Code}

The forward and inverse extracted procedures each contain three nested loops.
The invariants must express both control-state facts and semantic progress.

Typical outer-loop ingredients are:
\[
0 < \texttt{len} \le 128,
\qquad
\texttt{len} = 2^\ell,
\qquad
\code{poly\_repr}(\texttt{rp}, r),
\]
together with the statement that all fully completed stages of the transform
already agree on both sides.

Typical inner-loop ingredients are:
\[
\texttt{start} \le \texttt{j} \le \texttt{start} + \texttt{len},
\]
and equality of all coefficients outside the active butterfly window.

This task turns the local butterfly formulas into full procedure-level
refinement theorems.

\section{Task XII: Imperative-to-Abstract Identification}

After the extracted code has been related to the imperative model, one must
still prove the mathematical identities
\[
\texttt{NTT.ntt}(p,\zeta) = \mathrm{Rq.ntt}(p),
\qquad
\texttt{NTT.invntt}(p,\zeta^{-1}) = \mathrm{Rq.invntt}(p).
\]

The formulas governing this task are not low-level machine identities, but
rather structural equalities between:
\begin{itemize}
\item staged butterfly decompositions;
\item the bit-reversal exponents \(2\operatorname{br}(i)+1\);
\item the closed-form sums in \texttt{Rq.ec}.
\end{itemize}

This is where the Cooley--Tukey factorization must be matched to the explicit
ring-theoretic definition.

\section{Task XIII: Final Composition}

At the end of the project the formulas from all prior tasks are composed:
\[
\text{concrete state}
\xrightarrow{\code{poly\_repr}}
\text{abstract polynomial}
\xrightarrow{\mathrm{Rq.ntt}\ \text{or}\ \mathrm{Rq.invntt}}
\text{target polynomial},
\]
while the program side follows
\[
\texttt{poly\_ntt\_jazz}
\quad\text{or}\quad
\texttt{poly\_invntt\_jazz}.
\]

The completed chain should read:
\[
\texttt{poly\_ntt\_jazz}
\sim
\texttt{M.\_poly\_ntt}
\sim
\texttt{NTT.ntt}
=
\mathrm{Rq.ntt},
\]
and similarly
\[
\texttt{poly\_invntt\_jazz}
\sim
\texttt{M.\_poly\_invntt}
\sim
\texttt{NTT.invntt}
=
\mathrm{Rq.invntt}.
\]

\section{Task-to-Formula Table}

\begin{center}
\begin{longtable}{p{0.19\linewidth} p{0.31\linewidth} p{0.38\linewidth}}
\toprule
Task & Governing formula & Role in the proof \\
\midrule
\endfirsthead
\toprule
Task & Governing formula & Role in the proof \\
\midrule
\endhead
Imperative specs & forward and inverse butterfly recurrences & semantic target for extracted-code refinement \\
Losslessness & explicit integer variants for all loops & well-formedness of Hoare and pRHL reasoning \\
Root structure & \(zroot^{256}=-1\), \(zroot^{512}=1\) & algebraic foundation for inverse correctness \\
Representation & \(\code{poly\_repr}(\texttt{rp},p)\) & concrete-to-abstract state relation \\
Forward zetas & \(\texttt{jzetas}[k] \leftrightsquigarrow zroot^{\operatorname{br}(k)}R\) & table-read interpretation in forward butterflies \\
Inverse zetas & \(\texttt{jzetas\_inv}[k] \leftrightsquigarrow zroot^{-(\operatorname{br}(k)+1)}R\) & table-read interpretation in inverse butterflies \\
Inverse scale & \(\operatorname{incoeff}(-29720)R^{-1} = \operatorname{inv}(256)R\) & final inverse normalization \\
Montgomery multiply & \(a\cdot b\cdot R^{-1}\) modulo \(q\) & arithmetic core of extracted multiplication \\
Forward local step & \((x,y)\mapsto (x+\zeta y,\ x-\zeta y)\) & atomic forward simulation \\
Inverse local step & \((x,y)\mapsto (x+y,\ \zeta(x-y))\) & atomic inverse simulation \\
Loop refinement & invariants on \(\texttt{len},\texttt{start},\texttt{j},\texttt{zetasctr}\) & global simulation of extracted routines \\
Abstract identification & imperative staged transform \(=\) sum formula in \texttt{Rq.ec} & connection to mathematical specification \\
Final theorem & \(\{\code{poly\_repr}\}\;P\;\{\code{poly\_repr}\circ \mathrm{Rq.ntt}\}\) & end-to-end program correctness \\
\bottomrule
\end{longtable}
\end{center}

\section*{Conclusion}

The project admits a clean mathematical decomposition. Each task is governed by
a small family of formulas: arithmetic identities for roots of unity and
Montgomery constants; representation formulas linking words to coefficients;
local butterfly recurrences; global loop invariants; and finally the closed
transform formulas from the ring-theoretic specification. Writing the project
in this task-by-formula form makes its architecture transparent: the entire
verification is the controlled composition of these formulas inside the proof
language of EasyCrypt.

\end{document}
